\scanpage{18}
à des $f$ de type fini au lieu de $f$ propre. Se ramène au cas d'une \uline{imm.} \uline{ouverte}. Se résoud par l'affirmative lorsqu'on dispose de la résolution des singularités \inserted{$+$ pureté}, p.ex. en car. $0$ \inserted{pour ``bons \unsure{préschémas|premiers}''}, \unsure{ou plus} \inserted{bons \unsure{présch.}} \unsure{égale car.} lorsqu'on dispose de résolution des singularités \deleted{\unsure{en}} dans ce cadre....

14) \uline{Les} $R^i_!f_*(F)$.

Soit $f:X\to Y$ morphisme compactifiable, i.e. qui se factorise en
\diagram{\begin{tikzcd}
X \arrow[r,no head,"i","\text{imm. ouv.}"']&\overline X\arrow[r,"\overline f","\text{propre}"']&Y
\end{tikzcd}}
(p.ex. $f$ qu. proj. et $Y$ quasi-cpt).

Pour tt faisceau de torsion $F$ sur $X$, considérons $i_!(F)$ [exact en $F$] puis les $R^q_{\unsure{\deleted{!}|\xi}}\overline f_*(i_!(F))$, forment un foncteur coh. en $F$. Il résulte de façon quasi-triviale du th. de changement de base que ce foncteur coh.
\mspage{17}
